\iflanguage{english}
{\chapter{Einleitung}}
{\chapter{Introduction}}

\label{sec:introduction}


\section{Motivation}
3D human body reconstruction from a single image has emerged as a crucial challenge in computer graphics and computer vision due to its broad applicability in practical scenarios, such as \ac{VTO}, game industry, \ac{AR}, and artistic creation. Recently, the advances in neural rendering and generative models such as NeRF\cite{mildenhall2020nerf}, 3D Gaussian Splatting\cite{kerbl3Dgaussians} and Stable Diffusion\cite{rombach2021highresolution} have improved the ability to estimate the geometry and pose of human bodies from images. However, the accurate estimation of high-resolution surface UV textures remains a challenging problem, hindering the realism, consistency, and visual quality of reconstructed models，because the reconstruction quality of UV texture maps is very important to the final visual performance.

Previous pioneering work such as SMPLitex~\cite{casas2023smplitex} used diffusion models for texture inpainting from partial to full texture, they made it possible to estimate plausible and visually appealing textures from just a single image. Nevertheless, despite these promising results, such methods could only generate textures at relatively low resolutions (512×512 pixels), limiting the visual fidelity and detail required by demanding downstream applications. Moreover, SMPLitex often relies on limited training data, reducing the generalization capability to novel subjects and scenarios. Another recent work, TexDreamer~\cite{liu2024texdreamer} could estimate higher resolution human UV texture maps, but TexDreamer models training required more training data, which means more computing resources. 

Given these limitations, there is a pressing need to enhance existing methods to produce more realistic, detailed, and high-resolution human textures but relatively less resource-intensive. Leveraging \ac{SOTA} \ac{SR} models and advanced generative diffusion frameworks can potentially overcome these barriers, providing richer details and more robust solutions for texture estimation.

\section{Goal}
In order to address the previously mentioned challenges, this thesis proposes a novel method: High-Resolution Human Texture estimation (HRHumanTex), which is specifically designed to generate high-quality, high-resolution (1024×1024 pixels) and realistic UV textures from a single image of a clothed human. Specifically, this proposed work is centered around the following main goals:

\paragraph{Enhanced UV Data Generation}
Apply multiple \ac{SOTA} \ac{SR} techniques such as RealESRGAN~\cite{wang2021realesrgan}, SRFormer~\cite{zhou2023srformer}, DiffBIR~\cite{lin2024diffbir} and so on to elevate the original 512×512 SMPLitex texture dataset to 1024×1024 resolution, this thesis systematically benchmarks their strengths and weaknesses in enhancing UV-space textures.

\paragraph{High-Resolution Texture Completion}
With the improved datasets, this work trained multiple specialized inpainting models using the DreamBooth framework fine-tuning~\cite{ruiz2023dreamboothfinetuningtexttoimage} based on Stable Diffusion~\cite{rombach2021highresolution}. Each model is trained on textures enhanced by different SR methods, allowing me to explore how the quality of training data influences inpainting performance.

\paragraph{Systematic Evaluation Methodology}
To fairly evaluate the outcomes, this work established a robust evaluation pipeline combining traditional super-resolution metrics (e.g., \ac{PSNR}, \ac{SSIM}, \ac{LPIPS}, \ac{FID}) to pseudo-ground truth and rendering-based evaluation and optimization through differentiable rendering frameworks~\cite{ravi2020pytorch3d} to rendered pseudo-ground truth and real images.

\paragraph{Construction of an Integrated High Resolution Pipeline}
Finally, this work integrated all components--DensePose-based pixel-to-surface estimation, segmentation using SGHM, SR-enhanced data preparation, and fine-tuned diffusion-based inpainting models--into a unified pipeline that can produce high-quality 1024*1024 UV textures from a single image input.

By combining high-resolution super-resolution techniques, inpainting models, and thorough evaluation strategies, HRHumanTex seeks to push the boundaries of texture estimation. The resulting textures not only look more realistic but also hold greater value for downstream applications in virtual human creation, the game industry, and so on.

\section{Structure of the Work}

This thesis is structured into the following chapters:

Chapter 1: Introduction

This chapter outlines the motivation and research goals behind this work, with a distince focus on the importance of estimating high-resolution human textures from a single image. It also provides a blueprint of the entire thesis structure.

Chapter 2: Background

This chapter covers fundamental concepts necessary to understand this thesis, including 3D human reconstruction, UV texture mapping, and Stable Diffusion models~\cite{rombach2021highresolution}. It also covers fine-tuning techniques like  DreamBooth~\cite{ruiz2023dreamboothfinetuningtexttoimage}, human parsing tools such as DensePose~\cite{wu2019detectron2}, and \ac{SGHM} ~\cite{chen2022sghm}.

Chapter 3: Related Work

This chapter reviews existing literature in two main areas:

(1) 3D Human Reconstruction Review, comparing SMPL~\cite{SMPL:2015}-based explicit reconstruction methods and implicit reconstruction techniques like PIFu~\cite{saito2019pifu} and PIFuHD~\cite{saito2020pifuhd}.

(2) Human Texture Estimation Review, discussing approaches such as SMPLitex~\cite{casas2023smplitex}, TexFormer~\cite{xu2021texformer}, HPBTT~\cite{zhao2020human}, 360-degree Human Texture~\cite{lazova2019360}, and recently TexDreamer~\cite{liu2024texdreamer}. This chapter highlights their respective contributions and limitations and identifies their gaps with this thesis.

Chapter 4: Implementation

This chapter presents the full pipeline of the proposed HRHumanTex framework, including

Data enhancement stage using various \ac{SOTA} \ac{SR} models (RealESRGAN, DiffBIR, SRFormer, SwinIR, BSRGAN, StableSR, RCAN) and \ac{SR} evaluation strategy.

Training procedures for multiple texture-specialized DreamBooth inpainting models based on Stable Diffusion.

Integration of DensePose and SGHM for partial texture extraction, which are then completed into full-body texture maps through diffusion-based inpainting.

Ablation studies on data scale influences on model training and training hyperparameter configurations are also provided.

Experiments on differentiable rendering-based texture quality evaluation on optimization.

Chapter 5: Results

Presents both quantitative and qualitative experimental results of the proposed method, especially with the final rendering-based evaluation results and detailed analyses.

Chapter 6: Conclusion and Future Work

The final chapter summarizes the thesis's main contributions, discusses current limitations, highlights key findings, and includes possible improvements for future research. Including applications in virtual human modeling and interactive environments.

References:

A complete list of cited works, including DOIs, arXiv and github links identifiers where applicable, is provided for easy access to source materials

Appendices: 

Appendix A: Additional qualitative results and visual comparisons.

Appendix B: Selected special cases and failure cases with discussion.

Appendix C: Supplementary implementation notes and practical tips.

Appendix D: Acknowlegement.